\documentclass[sigplan,10pt,anonymous,review]{acmart}

\setcopyright{none}
\settopmatter{printacmref=false, printccs=false, printfolios=true}
\renewcommand\footnotetextcopyrightpermission[1]{}
\pagestyle{plain}

\acmConference[PAgE '27]{Principles of Agentic Engineering}{January 2027}{Mexico City, Mexico}
\acmBooktitle{Principles of Agentic Engineering (PAgE '27), January 2027, Mexico City, Mexico}

\begin{document}

\title{From Prompt Assembly to Governed State: Explicit Memory Boundaries for Tool-Using Agents}
\subtitle{Work in Progress}

\author{Anonymous Author(s)}
\affiliation{\institution{Anonymous Institution}}

\begin{abstract}
Tool-using language agents accumulate several kinds of state that are often
collapsed into one prompt: dialogue, tool observations, scratch data,
long-lived knowledge, authorization context, and the available action surface.
This conflation makes failures difficult to localize.  Truncation can separate
a tool invocation from its result; a retrieved fact can outlive its evidence;
and a memory query can cross a tenant boundary before the model ever sees the
data.  We present a work-in-progress architecture that treats agent context as
a governed state-assembly problem.  The architecture separates immutable
instructions, committed conversation, protected current input, spillable
working memory, and a provenance-bearing entity--fact store.  Context
capabilities are plugins sharing one tool manager, while compaction is a
single, observable pre-call transition that preserves tool-call/result pairs.
Persistent memory enforces ownership and principal-based visibility at the
storage boundary and represents non-aggregate updates by supersession rather
than destructive replacement.  We formulate seven runtime invariants covering
budget accounting, transcript integrity, capability coherence, provenance,
and authorization.  We then propose a fault-injection and long-horizon
evaluation measuring task success, grounded recall, unauthorized disclosure,
token cost, and latency against rolling-window, summarization, and retrieval
baselines.  The implementation is an open-source, strict-TypeScript agent
runtime; its identity and artifact URL are withheld during anonymous review.
The goal is not to claim that memory makes an agent safe, but to show how
explicit state boundaries turn important safety assumptions into interfaces,
checks, and testable hypotheses.
\end{abstract}

\keywords{language agents, context management, memory, tool use, runtime
invariants, access control}

\maketitle

\section{Introduction}

Language agents interleave model inference with external actions.  Systems
such as ReAct make the reasoning--action loop explicit~\cite{yao2023react},
while Toolformer demonstrates that models can learn when and how to call
APIs~\cite{schick2023toolformer}.  Once such agents operate for more than a
few steps, however, their engineering problem is no longer only ``what should
the model do next?''  It is also ``which state should be visible, mutable,
retrievable, and authorized at this step?''

Existing architectures have shown the value of modular memory
components~\cite{sumers2023coala}, dynamic retrieval and
reflection~\cite{park2023generative}, and operating-system-inspired movement
between memory tiers~\cite{packer2023memgpt}.  We build on this direction but
focus on a complementary question: which invariants should an agent runtime
enforce even when model behavior is unreliable?  The distinction matters
because many failures occur outside the model.  A context manager may orphan a
tool result during truncation, retrieve a superseded fact, expose a record to
the wrong principal, or maintain two divergent registries of executable tools.
More capable reasoning does not repair these faults.

This paper reports work in progress on a strict-TypeScript runtime for
tool-using agents.  Its design separates five state classes and gives each a
different lifetime and mutation policy.  A plugin protocol contributes prompt
content and tools to a single assembly pipeline.  A context budget is computed
before each model call, and compaction occurs once at that boundary.  Large or
old tool results may move to working memory, but historical tool invocations
and results are removed atomically.  Long-term memory is a typed entity--fact
system with confidence, importance, evidence, provenance, temporal validity,
and supersession.  Authorization is evaluated before retrieved records become
model input.

The contributions of this work-in-progress paper are:

\begin{itemize}
  \item a state taxonomy and preparation semantics for assembling bounded
        model context from independently governed components;
  \item seven executable invariants spanning context integrity, capability
        coherence, persistent-memory provenance, and access control;
  \item a reference implementation in a vendor-neutral agent runtime; and
  \item an evaluation design combining property tests, fault injection, and
        long-horizon agent tasks, including explicit safety and cost metrics.
\end{itemize}

We deliberately do not report end-task improvements yet.  The current evidence
is architectural and test-based; the proposed controlled evaluation is the
next stage.  This distinction is important for avoiding the common error of
equating the existence of a memory subsystem with improved reliability.

\section{Failure Model and Design Requirements}

We assume a model that can produce incorrect text, malformed tool arguments,
unwarranted memory writes, and inconsistent plans.  External tool results can
be large, stale, adversarial, or partially trusted.  The host application,
rather than the model, authenticates users and resolves organizational
membership.  Storage and runtime code are inside the trusted computing base;
model output and retrieved external content are not.

This model exposes four recurring failure classes.

\paragraph{Context-integrity failures.}
A token-pressure policy removes a tool invocation but retains its result, or
vice versa.  The next provider request is malformed or, worse, semantically
ambiguous.  Repeated compaction at several layers can also summarize already
summarized evidence and make loss difficult to audit.

\paragraph{State-lifetime failures.}
Ephemeral scratch data is promoted to durable user knowledge, or long-lived
facts exist only in a rolling prompt.  A corrected fact destructively
overwrites its predecessor, erasing provenance and making historical queries
impossible.

\paragraph{Capability-coherence failures.}
The prompt describes one tool set while the executor uses another.  A context
plugin maintains its own registry, so disabling a capability removes its
instructions but leaves an executable function reachable.

\paragraph{Authorization failures.}
A model-generated identity or group is treated as authoritative.  Filtering
occurs after retrieval, allowing unauthorized data into prompt construction,
logs, caches, or ranking.  AgentDojo's results show why untrusted tool content
must be treated as an active security input rather than inert
text~\cite{debenedetti2024agentdojo}.

From these failures we derive four requirements: state classes must have
explicit lifetimes; every model request must be assembled at one observable
boundary; tool capability must have a single source of truth; and memory
authorization must occur before model visibility.

\section{Architecture}

\subsection{State classes}

At turn $t$, the runtime state is

\begin{equation}
S_t = \langle P, H_t, U_t, W_t, K_t, M_t \rangle ,
\end{equation}

where $P$ contains the system contract and enabled-plugin instructions;
$H_t$ is committed conversation history; $U_t$ is the current input; $W_t$
is spillable working memory; $K_t$ is persistent entity--fact knowledge; and
$M_t$ is the single tool manager.  Table~\ref{tab:state} summarizes their
governance.

\begin{table*}[t]
\caption{State classes are separated by lifetime, authority, and loss policy.}
\label{tab:state}
\centering
\small
\begin{tabular}{lllll}
\toprule
State & Typical contents & Lifetime & Writer & Pressure behavior \\
\midrule
$P$: contract & system prompt, plugin rules & agent configuration & host/plugins & never compacted \\
$H_t$: history & user/assistant messages, tool pairs & session & runtime & pair-aware eviction \\
$U_t$: current input & new request or newest tool results & one model step & host/runtime & protected; reject or truncate explicitly \\
$W_t$: working memory & large observations, scratch artifacts & session or task & agent/runtime & indexed spill and explicit retrieval \\
$K_t$: knowledge & entities, facts, evidence, temporal state & cross-session & host/ingestor/agent policy & ranked, scoped retrieval \\
$M_t$: capabilities & executable tool definitions and policy & agent lifetime & host/plugins & definitions budgeted, never silently dropped \\
\bottomrule
\end{tabular}
\end{table*}

This separation does not imply that every application should enable every
class.  Features are plugins: a disabled feature contributes neither prompt
content nor tools.  Retrieval-only deployments can enable $K_t$ reads while
disallowing model-authored writes.  Session ingestors can observe a transcript
before compaction and update memory asynchronously, without blocking the model
call.

\subsection{One preparation boundary}

The preparation function computes

\begin{equation}
Q_t, B_t = \operatorname{prepare}(S_t, L, R),
\end{equation}

where $L$ is the model context limit, $R$ is the response reserve, $Q_t$ is the
provider-ready sequence, and $B_t$ is an inspectable budget.  The budget
contains separate terms for tool definitions, system and plugin content,
conversation, and current input:

\begin{equation}
b_P + b_H + b_U + b_M \leq L - R.
\end{equation}

Tool schemas and protected current input are accounted for before historical
conversation.  If the system contract plus current input cannot fit, the
runtime follows an explicit oversized-input policy rather than silently
evicting a different state class.  Compaction is triggered only at a configured
critical utilization threshold and runs once immediately before the provider
call.  Budget and compaction events are emitted for monitoring.

The implemented algorithmic policy first identifies historical tool pairs by
call identifier.  Large results are placed in $W_t$ with a descriptive key;
old pairs beyond a configured cap can be spilled; and a rolling window removes
remaining old messages.  Every removal operation treats a call and its result
as one unit.  Immediately before dispatch, a sanitizer rejects or repairs
unmatched historical pairs.  This is intentionally a runtime property rather
than a prompting convention.

\subsection{Plugin-scoped context, shared tools}

Each context plugin exposes four conceptually distinct surfaces: static
instructions, dynamic model-visible content, serializable state, and optional
tools.  Plugins track the token size of their own dynamic content and declare
whether it can be compacted.  The context manager assembles all enabled plugin
instructions and content into one developer message.

All tools, including plugin-provided and late-bound tools, are registered in
the same manager $M_t$.  Consequently, tool visibility, enablement, permission
checks, and execution observe the same state.  This removes a class of
time-of-check/time-of-use mismatches between prompt assembly and execution.
Dynamic tool catalogs can load a bounded subset of a large tool collection,
but the loaded subset still enters $M_t$ and the token budget.

\subsection{Persistent knowledge with provenance}

The persistent store represents identity as entities and knowledge as facts.
An atomic fact is approximately a tuple

\begin{equation}
f = \langle s,p,o,c,i,e,z,[v_0,v_1],\sigma \rangle,
\end{equation}

where $s,p,o$ are subject, predicate, and object/value; $c$ is confidence;
$i$ is importance; $e$ is supporting evidence; $z$ identifies the source
signal; $[v_0,v_1]$ is a validity interval; and $\sigma$ optionally identifies
a superseded predecessor.  Facts can additionally bind to several contextual
entities.  Document facts store longer narratives such as profiles.

Non-aggregate corrections create a successor and archive the prior fact rather
than mutating history in place.  Queries can therefore apply an ``as of'' time
and exclude superseded knowledge.  Retrieval ranks visible candidates using
confidence, recency, predicate weight, and importance; semantic retrieval and
bounded graph traversal are optional tiers.  Entity resolution proceeds from
strong identifiers through normalized names and aliases to a confidence-capped
semantic fallback.  The model can emit surface forms, while trusted runtime
code resolves stored identifiers.

This model is intentionally more structured than storing transcript chunks in
a vector index.  The vector representation is a retrieval aid; provenance,
temporal state, ownership, and graph relationships remain first-class fields.

\subsection{Authorization before model visibility}

Every entity and fact has an owner.  The host supplies a trusted set of caller
principals (for example user, group, entity, and service identities).  At the
storage boundary, record permissions and explicit ACL grants are materialized
as read and write principal sets.  A record $r$ is visible to caller $a$ only
if

\begin{equation}
\operatorname{Read}(r) \cap \operatorname{Principals}(a) \neq \emptyset .
\end{equation}

The query performs this intersection; it is not a post-retrieval model filter.
Write principals are constrained to be a subset of read principals.  An empty
caller-principal set fails closed.  Because pre-migration records lack the
materialized sets, principal mode has no legacy fallback: a backfill is an
explicit deployment gate.  The host remains responsible for constructing the
principal set from authenticated state.  Allowing model output to add a group
principal would move the vulnerability outside the library rather than solve
it.

\section{Executable Invariants}

We use invariants both as design documentation and as property-test targets.
They do not prove end-to-end agent correctness; they narrow the set of failures
that the stochastic component can induce.

\begin{description}
  \item[I1: Budget accountability.] Every provider request exposes a budget
  decomposition, and fixed content plus response reserve must fit before
  historical content is admitted.

  \item[I2: Current-input preservation.] Accepted $U_t$ is present in $Q_t$;
  ordinary history compaction cannot remove the newest request or newest tool
  observation.

  \item[I3: Historical pair closure.] For each historical call identifier in
  $Q_t$, its tool invocation and result are either both present or both absent.

  \item[I4: Single preparation transition.] Compaction occurs at most once per
  provider call and is observable through a transition log and budget event.

  \item[I5: Capability coherence.] The agent executor and context assembler
  reference the same manager: $M_t^{agent} \equiv M_t^{context}$.  Disabling a
  plugin removes its contributed content and tools.

  \item[I6: Provenance-preserving correction.] A non-aggregate correction does
  not destroy its predecessor; the successor links to it and temporal queries
  select the appropriate visible version.

  \item[I7: Authorized retrieval.] Every returned record is readable by at
  least one trusted caller principal, and write principals are a subset of read
  principals.
\end{description}

The current implementation has unit and integration tests for pair-preserving
compaction, context snapshots and persistence, plugin feature flags, shared
tool-manager identity, owner and ACL matrices, supersession, temporal queries,
entity resolution, and fail-closed principal behavior.  These tests establish
implementation conformance to local properties.  They do not yet establish
whether agents complete more tasks or whether the ranking policy selects the
best evidence under realistic load.

\section{Evaluation Plan}

Our evaluation separates deterministic runtime properties from stochastic
agent outcomes.

\subsection{Research questions}

\textbf{RQ1: Integrity.} Under token pressure and malformed histories, does
pair-aware preparation eliminate structurally invalid provider requests while
preserving the current observation?

\textbf{RQ2: Long-horizon reliability.} Do explicit working and persistent
memory tiers improve task completion and evidence-grounded answers relative to
rolling-window and summarization baselines at equal model context limits?

\textbf{RQ3: Authorization.} Under cross-tenant queries, stale indexes, entity
merges, and adversarial tool content, does storage-boundary filtering prevent
unauthorized records from entering model-visible context?

\textbf{RQ4: Tradeoffs.} What task-success gains, if any, justify additional
tokens, retrieval operations, and p50/p95 latency?

\subsection{Workloads and baselines}

We plan three workload families.  First, repository-navigation tasks require
an agent to inspect many files and retain tool observations across a long
trajectory; a held-out subset of SWE-bench-style issues provides realistic
software-engineering structure~\cite{jimenez2024swebench}.  Second, synthetic
assistant tasks combine email, calendar, people, projects, corrections, and
time-sensitive questions.  Their generated ground-truth graph makes
provenance, staleness, and authorization directly measurable.  Third,
tool-policy dialogues adapt the state-based evaluation style of
$\tau$-bench~\cite{yao2024taubench}, including repeated trials to measure
reliability rather than best-case success.

We will compare: (B0) full history until rejection; (B1) oldest-first rolling
window; (B2) rolling summaries; (B3) embedding retrieval over unstructured
chunks; (B4) tiered context without persistent knowledge; and (B5) the full
architecture.  All conditions use the same model, tools, response reserve,
and nominal context limit.  Retrieval-only and model-write-enabled memory are
reported separately because they have different safety envelopes.

\subsection{Fault injection}

Deterministic tests inject orphaned tool calls/results, oversized Unicode tool
outputs, repeated compaction thresholds, plugin failures, serialization and
restore cycles, conflicting facts, expired facts, low-confidence entity
matches, revoked ACL entries, incomplete principal sets, pre-backfill records,
and entity merges.  Security cases place prompt-injection strings inside tool
results and memory evidence.  We do not expect state separation alone to stop
all prompt injection; the question is whether authorization and transcript
invariants remain intact while task-level defenses are evaluated independently.

\subsection{Metrics and analysis}

Primary safety metrics are malformed-request rate, current-input loss,
unauthorized-record exposure (target: zero), and provenance coverage of facts
used in final answers.  Reliability metrics are state-goal success,
evidence-supported answer accuracy, stale-fact selection, and repeated-trial
$\mathrm{pass}^k$.  Systems metrics are input/output tokens, peak prompt
utilization, memory reads/writes, compaction count, and p50/p95 latency.

Each stochastic condition will run paired task instances across multiple
seeds.  We will report confidence intervals and per-task deltas, not only
aggregate means.  Safety failures will be listed individually.  Model and
runtime versions, prompts, synthetic data generators, and event traces will be
pinned in the artifact.  The principal hypothesis is that explicit tiers and
pair-aware compaction reduce structural failures and improve long-horizon
success at a measurable latency cost.  A null or negative task-success result
would remain useful if the invariant and cost results identify where the
architecture does and does not help.

\section{Discussion}

\paragraph{Boundaries are not policies.}
The architecture makes state classes and enforcement points explicit, but a
host can still choose an unsafe visibility default, trust an incorrect
principal set, or permit the model to write unsupported facts.  The system
therefore exposes retrieval-only operation, configurable visibility policy,
evidence requirements, and confidence thresholds.  Defaults and deployment
guidance are part of the safety surface.

\paragraph{Compaction is lossy computation.}
Moving a tool result to working memory preserves bytes but changes immediate
availability; summarizing it changes content.  Both operations should be
treated as semantic transitions with logs, not invisible prompt optimization.
Our evaluation will distinguish retained, retrievable, and irrecoverably
dropped observations.

\paragraph{Structured memory introduces its own errors.}
Entity resolution can merge distinct people, extraction can invent a fact, and
ranking can prefer salient but stale evidence.  Provenance and supersession
make these errors inspectable but do not prevent them.  Restraint-oriented
ingestion---requiring evidence quotes and using a skeptical second pass---is an
available policy whose effect still requires measurement.

\paragraph{Generality.}
The concepts are provider-neutral, but the reference implementation and tests
are TypeScript-specific.  Provider message schemas also differ.  Pair closure
is expressed over a normalized internal representation, and adapters must
preserve call identifiers correctly.

\section{Agent-Assisted Research Workflow}

The workshop explicitly invites descriptions of agent use in research.  An AI
coding agent was used to inspect the implementation and documentation, locate
relevant tests, search for primary literature, and produce the initial paper
scaffold and prose.  Human authors selected the research claim, constrained
the paper to work-in-progress status, and are responsible for validating every
architectural statement, citation, experiment, and anonymization decision.
No empirical result in this draft was generated or inferred by the agent.

For reproducibility, the final artifact will record the agent instructions,
the source revision inspected, generated diffs, and the commands used for
tests and paper compilation.  Agent-written text is treated as an untrusted
draft: claims must be matched to code or measured data, references must resolve
to primary sources, and proposed experiments require human approval before
execution.  The main lesson so far is that an agent can rapidly connect a broad
implementation to a paper structure, but it tends to turn implemented
mechanisms into implied benefits.  Explicitly separating ``implemented,''
``tested invariant,'' and ``measured outcome'' is a useful oversight rule.

\section{Conclusion}

Safe agentic engineering requires more than improving the model's next-action
prediction.  The surrounding runtime decides what the model can see, remember,
forget, and execute.  We have presented a work-in-progress architecture that
separates prompt contract, conversation, current input, working memory,
persistent knowledge, and tool capability; assembles them at one budgeted
boundary; and enforces integrity and authorization invariants outside the
model.  The immediate next step is a controlled evaluation that tests whether
these structural guarantees translate into grounded long-horizon behavior at
an acceptable cost.  Regardless of outcome, stating the invariants makes the
system's safety claims falsifiable.

\bibliographystyle{ACM-Reference-Format}
\bibliography{references}

\end{document}
